\section{Funzioni, Scope e Stack (Lez. 06-07)}

\subsection{Cosa sono le Funzioni}
Una funzione è un blocco di codice progettato per eseguire un compito specifico. La esegui "chiamandola" (invocandola).
Serve a non dover riscrivere lo stesso codice tante volte.

\begin{concept}{Definizione Funzione}
Una funzione è un blocco di codice progettato per eseguire un compito specifico. La esegui "chiamandola" (invocandola).
\end{concept}
\begin{jsbox}
// Definizione
function saluta(nome) {
    return "Ciao " + nome + "!";
}

// Utilizzo
let messaggio = saluta("Mario"); 
console.log(messaggio); // Stampa: "Ciao Mario!"
\end{jsbox}

\subsection{Scope: Dove vivono le variabili}
In JavaScript, le variabili non sono visibili ovunque. Dipende da dove le dichiari.

\subsubsection{1. Block Scope (let, const)}
\begin{concept}{Scope di Blocco (Recinto Privato)}
Le variabili definite con \texttt{let} o \texttt{const} dentro un blocco \texttt{\{ \}} (come un if o un for) esistono \textbf{solo lì dentro}.
\end{concept}
\begin{jsbox}
// Scope limitato (dentro un blocco if)
if (true) {
    let x = 2;
}
// x qui non esiste piu'
\end{jsbox}

\subsubsection{2. Function Scope (var)}
\begin{concept}{Scope di Funzione (Old School)}
Le variabili definite con \texttt{var} dentro una funzione sono visibili in tutta la funzione, ma non fuori.
\end{concept}
\begin{jsbox}
function test() {
    var car = "Fiat";
}
// car qui non esiste
\end{jsbox}

\subsection{Arrow Functions}
\begin{concept}{Sintassi Concisa}
Le Arrow Functions \texttt{() =>} semplificano la scrittura. Se c'è una sola espressione, il \texttt{return} è implicito.
\end{concept}
\begin{jsbox}
// Prima
let hello = function() {
    return "Hello World!";
}

// Dopo (Arrow Function)
let hello = () => {
    return "Hello World!";
}

// Se hai una sola riga, puoi togliere le graffe e il return:
let hello = () => "Hello World!";
\end{jsbox}

\subsection{Closures (Chiusure)}
\begin{concept}{Memoria Persistente}
Una closure è quando una funzione figlia "ricorda" le variabili dello scope padre anche dopo che il padre è terminato.
\end{concept}
\begin{jsbox}
function creaContatore() {
    let count = 0;
    return function() {
        count += 1; // Ricorda 'count'
        return count;
    }
}

const conta = creaContatore();
console.log(conta()); // 1
console.log(conta()); // 2
console.log(conta()); // 3
\end{jsbox}

\subsection{Ricorsione}
\begin{concept}{Funzione che chiama se stessa}
È utile per risolvere problemi che si possono spezzare in problemi più piccoli uguali (come il Countdown).
\end{concept}
\begin{jsbox}
function countdown(n) {
    if (n <= 0) {
        console.log("Via!");
    } else {
        console.log(n);
        countdown(n - 1); // Chiama se stessa con n-1
    }
}
countdown(3);
// Output: 3, 2, 1, Via!
\end{jsbox}
